% Truth-level surrogate models for machine-learned emerging-jet searches - draft
% JHEP preprint style (jheppub.sty). Build: pixi run paper
\documentclass[a4paper,11pt]{article}
\usepackage{jheppub}
\usepackage{amssymb}
\usepackage{xcolor}
\usepackage{subcaption}
\usepackage{booktabs}
\captionsetup[subfigure]{skip=1pt}
\setlength{\emergencystretch}{1em}

\graphicspath{{../results/figures/}}

% ------------------------------------------------------------------ macros
\newcommand{\stub}[1]{\begin{center}\fbox{\parbox{0.92\textwidth}{\color{red!70!black}\textbf{STUB:} #1}}\end{center}}
\newcommand{\todo}[1]{{\color{red!70!black}[TODO: #1]}}
\newcommand{\dnought}{d_0}
\newcommand{\znought}{z_0}
\newcommand{\ctau}{c\tau}
\newcommand{\pid}{\pi_{\mathrm{D}}}
\newcommand{\rhod}{\rho_{\mathrm{D}}}
\newcommand{\qD}{q_{\mathrm{D}}}
\newcommand{\dR}{\Delta R}
\newcommand{\pT}{p_{\mathrm{T}}}
\newcommand{\mpid}{m_{\pid}}
\newcommand{\effSR}{\varepsilon_{\mathrm{SR}}}
% unnumbered bold heading on its own line
\newcommand{\obshead}[1]{\par\medskip\noindent{\bfseries #1}\par\nobreak\smallskip}

% Alternative titles, to choose from:
%   Reinterpreting machine-learned LHC searches with truth-level surrogate models
%   Truth-level surrogates for detector-level machine-learned taggers
%   Learning the detector response of machine-learned taggers from generator truth
%   Surrogate models for the reinterpretation of low-level machine-learning selections at the LHC
%   From generator to tagger: truth-level surrogates for machine-learned searches
\title{Reinterpreting machine-learned LHC searches with truth-level surrogate models}
\author{Jackson Burzynski}
\emailAdd{jackson.carl.burzynski@cern.ch}

\abstract{The reinterpretation of published LHC results with new
theoretical models is one of the most efficient ways to extend the
enormous investment represented by a dedicated search, but an
increasing number of analyses rely on machine-learning-based selection
algorithms operating on low-level detector inputs, whose response to
new signals cannot be assessed without full detector simulation and
reconstruction. We study surrogate models trained on truth-particle
information to approximate the detector-level efficiency of such an
analysis. The method is generic: a selection built on a custom jet
tagger with constituent-level inputs is replaced by a network that
consumes the generator record of the same jet and predicts the tagger
decision, and per-jet predictions are combined into event-level
acceptances without further approximation. We demonstrate it in a
fully controlled setting with a search for emerging jets as the
example, chosen because its signal topology depends on many dark-sector
parameters and its selection rests on a transformer tagger over
reconstructed tracks. A tagger is trained on the tracks inside
large-radius jets, a signal region is defined by two tagged jets at a
working point of $10^{3}$ background rejection, and a second
transformer operating only on the generator-level particles inside a
jet is trained to predict the tagger decision. The surrogate
reproduces the detector-level signal-region efficiency to
$5$--$10\%$ on held-out samples of its training models, across
dark-pion proper lifetimes from $0.01$ to $5$~mm and a scan of the
dark-sector fragmentation, and generalizes to mediator masses,
dark-pion masses and dark couplings never seen in training to within
$30\%$ across most of that model space, at a computational cost that
bypasses detector simulation entirely. Nothing in the construction is specific to dark showers or to
jets, and the same strategy applies to any analysis whose selection is a
machine-learned function of low-level reconstructed objects, including
event-level classifiers.}

\begin{document}
\maketitle

% =====================================================================
\section{Introduction}
\label{sec:intro}

One of the primary objectives of the LHC is a comprehensive search
program for beyond the Standard Model (BSM) phenomena, including new
heavy resonances, dark matter, supersymmetric particles, and exotic
long-lived particles. Despite the breadth of searches performed by the
LHC experiments thus far, no statistically significant evidence of new
physics has been observed. A central challenge of the LHC program is
that the space of viable BSM scenarios is too broad to be fully covered
by dedicated searches. Physics analyses at the LHC represent an
enormous investment of time and expertise, with a typical search
taking more than two years and multiple PhD-student FTEs to complete.
The reinterpretation of published results using new theoretical models
is one of the most efficient ways to extend this investment. However,
processing such models through the full chain of detector simulation,
reconstruction, and collaboration-specific analysis pipelines is
computationally prohibitive given the breadth of potential BSM models.
Moreover, only a small fraction of LHC analyses provide the input
material needed for external groups to carry out reinterpretation
using particle-level information alone. Existing reinterpretation
tools such as \textsc{Contur}~\cite{Butterworth:2016sqg} have made
important progress by leveraging the measurements preserved in the
\textsc{Rivet} library~\cite{Buckley:2010ar} to test BSM scenarios
without requiring detector simulation or experiment-specific
reconstruction algorithms. By construction, however, such tools are
limited to unfolded fiducial measurements and cannot assess the
sensitivity of dedicated BSM searches. This limitation is becoming
increasingly significant as more analyses rely heavily on
machine-learning-based selection algorithms operating on low-level
detector inputs, whose response to new theoretical models cannot be
easily assessed with existing fast-simulation frameworks such as
\textsc{Delphes}~\cite{deFavereau:2013fsa}.

One promising approach to addressing this challenge is the use of a
\emph{surrogate model} trained on truth-particle-level information to
approximate the detector-level efficiency of a given analysis,
bypassing the need for full detector simulation and reconstruction.
The surrogate model concept has been demonstrated in
Ref.~\cite{Corpe:2025rtz}, which trained a boosted decision tree to
approximate the selection efficiency of a search for displaced jets in
the ATLAS calorimeter~\cite{ATLAS:CalRatio}. That analysis, however,
selects on a small number of high-level per-jet quantities. The open
question addressed here is whether the same strategy survives contact
with a selection built on a low-level machine-learned tagger, in
which the detector response to a new signal is encoded in the
correlations among tens of reconstructed tracks per jet rather than in
a handful of physically motivated variables.

The strategy studied here is generic. Any selection built on a
machine-learned tagger with low-level inputs defines, for each object it
acts on, a binary detector-level outcome. A surrogate is a network that
consumes the generator-level description of the same object and is
trained to predict that outcome, so that the per-object probabilities it
returns can be combined into the event-level acceptance of the analysis
from the truth record alone. Nothing in this construction depends on
the physics of the signal, on the architecture of the tagger, or on the
tagger operating on jets: a classifier acting on the full event, with
the surrogate consuming the full generator record, is the same problem
with a larger input set.

Emerging jets~\cite{Schwaller:2015gea} serve as the example throughout
this paper. They are a signature of strongly interacting dark
sectors~\cite{Strassler:2006im,Han:2007ae,Bai:2013xga,Albouy:2022cin}
in which long-lived dark particles produced inside a jet decay back to
Standard Model particles at displaced vertices, and they are a
demanding test case for two reasons. These dark-QCD scenarios are
governed by many free parameters, including the strength of
interactions between dark particles, their masses, and their lifetimes,
and small changes in these parameters produce non-trivial shifts in the
final-state topology, so the sensitivity of an analysis cannot be
assumed to generalize across the parameter space without explicit
reinterpretation. The complex signal topology, with an arbitrary number
of dark particles per jet each decaying at lifetime-dependent
displacements, also makes it impractical to parameterize the selection
efficiency as a simple function of model parameters. Furthermore, the
current ATLAS search~\cite{ATLAS:EJ} relies on a transformer-based jet
tagger operating on low-level detector inputs whose response to new
signals cannot be evaluated without full detector simulation and
reconstruction. A surrogate that works here should work for simpler
signatures.

This paper builds and validates such a surrogate end to end in a
fully controlled setting. Signal and background events are generated
with \textsc{Pythia}~8 and passed through a \textsc{Delphes} simulation
of an ATLAS-like detector with realistic track impact-parameter
resolution. A transformer tagger trained on the reconstructed tracks
inside large-radius jets plays the role of the analysis, and a signal
region is defined by requiring two tagged jets. A second transformer,
operating only on the generator-level particles inside a jet, is then
trained to estimate the probability that the jet would be tagged by
the detector-level classifier. Because every element of the chain is
available, the surrogate can be confronted with the full
detector-level answer not only on the model it was trained on but on
dark-pion masses, lifetimes and couplings it has never seen, which is
precisely the situation in which a surrogate would be used.

Our goals are as follows.
(i) Establish that truth-level particle information is sufficient to
learn a meaningful approximation of a low-level track-based tagger.
(ii) Quantify the accuracy with which the surrogate reproduces the
event-level signal-region efficiency, both on held-out samples of the
training model and on models with different dark-sector parameters.
(iii) Identify which truth-level quantities carry the information the
surrogate uses, and where the approximation degrades.
(iv) Provide a concrete template for the material an analysis team
would publish to make a machine-learned search reinterpretable.

% =====================================================================
\section{Simulation and event reconstruction}
\label{sec:simulation}

\subsection{Signal and background samples}

Signal events are generated with \textsc{Pythia}~8.312
\cite{Bierlich:2022pfr} using the Hidden Valley
module~\cite{Carloni:2010tw}: $s$-channel production
$pp \to Z' \to \qD\bar{\qD}$ at $\sqrt{s} = 13.6$~TeV with
$m_{Z'} = 1.5$~TeV, followed by a dark shower and hadronization with
$N_{\mathrm{gauge}} = 3$, running dark coupling with
$\Lambda_{\mathrm{D}} = 10$~GeV, $m_{\qD} = 10$~GeV,
$\mpid = 5$~GeV, $m_{\rhod} = 20$~GeV, and vector fraction $0.75$,
following the unflavored benchmark of the Dark Showers Task
Force~\cite{Albouy:2022cin} and the original emerging-jets
proposal~\cite{Schwaller:2015gea}. Dark rhos decay promptly to dark
pion pairs, and dark pions decay to $d\bar d$ with proper lifetime
scanned over $\ctau(\pid) \in \{0.01, 0.05, 0.1, 0.5, 1, 5\}$~mm. This lifetime range
is deliberately short. At these lifetimes the displaced decay products
remain inside the pixel volume and the discrimination between dark
showers and QCD rests on impact parameters of a few tens of
micrometers and on substructure, which is the regime in which a
surrogate must reproduce the detector response most faithfully.

To test the generalization of the surrogate, two additional signal
models are generated in which the entire dark sector is rescaled
coherently with the dark-pion mass, $\Lambda_{\mathrm{D}}$, $m_{\qD}$,
$\mpid$, $m_{\rhod}$ and the shower cutoff all scaling with
$\mpid/5$~GeV, so that the shower remains self-consistent rather than
changing one mass in isolation. The alternative points have
$\mpid = 10$~GeV and $\mpid = 2$~GeV, with the same lifetime grid.
These samples are never used in training.

Beyond the mass rescaling, the fragmentation of the dark shower is
varied at fixed $\mpid$ and $\ctau$ by scanning the confinement scale
over $\mpid/\Lambda_{\mathrm{D}} \in \{0.2, 0.35, 0.5, 0.7, 1.0,
1.4\}$, with $m_{\qD} = \Lambda_{\mathrm{D}}$, the shower cutoff
following $\Lambda_{\mathrm{D}}$ and
$m_{\rhod} = \max(2\Lambda_{\mathrm{D}}, 2.2\,\mpid)$ so that the
forced $\rhod \to \pid\pid$ decay stays open, and by raising the
number of degenerate dark-quark flavours from one to two and three.
Two of the scan points, $\mpid/\Lambda_{\mathrm{D}} = 0.2$ and $1.0$,
are also generated with high statistics as statistically independent
samples and enter the training of the surrogate, while the remaining
points are reserved for evaluation.

The background is QCD dijet production with $\hat p_T > 450$~GeV,
matched to the $Z'$ decay spectrum. The same sequence of
hard-scattering configurations is reused across the lifetime grid, so
that the hard process is common to all lifetime points and the
lifetime comparison is a controlled one. The samples, their sizes and
their roles are collected in Table~\ref{tab:samples}. Each sample is
divided into statistically independent subsamples for training,
validation and testing, so that no event contributes to more than one
role, and the events that train the surrogate are disjoint from every
event the tagger has seen.

\begin{table}[htbp]
  \centering
  \small
  \setlength{\tabcolsep}{3pt}
  \begin{tabular}{llll}
    \toprule
    sample & $\ctau$ grid & tagger events (train / val / test) & surrogate training events \\
    \midrule
    QCD dijet & -- & $6.0\,/\,1.0\,/\,3.0 \times 10^{5}$ & -- \\
    nominal signal & 6 points & $14\,/\,1.8\,/\,4.2 \times 10^{5}$ & $1.2\times10^{6}$ \\
    $\mpid/\Lambda_{\mathrm{D}} = 0.2$ & 6 points & -- & $6.0\times10^{5}$ \\
    $\mpid/\Lambda_{\mathrm{D}} = 1.0$ & 6 points & -- & $6.0\times10^{5}$ \\
    \midrule
    $\mpid = 10$, $2$~GeV & 6 points & \multicolumn{2}{l}{evaluation only, $3.0\times10^{5}$ events each} \\
    $\mpid/\Lambda_{\mathrm{D}} = 0.35$, $0.7$, $1.4$ & $0.1$~mm & \multicolumn{2}{l}{evaluation only, $5.0\times10^{4}$ events each} \\
    $N_{\mathrm{flav}} = 2, 3$ & $0.1$~mm & \multicolumn{2}{l}{evaluation only, $5.0\times10^{4}$ events each} \\
    \bottomrule
  \end{tabular}
  \caption{Monte Carlo samples and their roles. The nominal signal
  is the dark-shower benchmark with $\mpid = 5$~GeV, the background is
  QCD dijet production with $\hat p_T > 450$~GeV, and the six-point
  lifetime grid is
  $\ctau(\pid) \in \{0.01, 0.05, 0.1, 0.5, 1, 5\}$~mm, with equal
  statistics per lifetime point. All subsamples are statistically
  independent. The training events of the surrogate, including those
  of the two $\mpid/\Lambda_{\mathrm{D}}$ training points, share no
  event with anything the tagger has seen or with any evaluation
  sample, and the two training points also have an independent
  $5.0\times10^{4}$-event evaluation sample at $\ctau = 0.1$~mm. The
  scan points not listed as training, as well as the alternative-mass
  models, are never used in training.}
  \label{tab:samples}
\end{table}

Figure~\ref{fig:samplespan} shows why the training set of the
surrogate is composed this way. The decision of the tagger correlates
with two properties of a jet that the model parameters steer almost
independently. The dark-pion lifetime sets how much of the jet is
displaced, and $\mpid/\Lambda_{\mathrm{D}}$ sets how the jet
fragments, that is, the multiplicity and hardness of the dark hadrons.
The lifetime grid varies the first at fixed fragmentation, and the two
$\mpid/\Lambda_{\mathrm{D}}$ training points vary the second at
every lifetime, so the surrogate is trained across the plane of the
key correlations the detector-level tagger has learned rather than
along a single line through it.

\begin{figure}[htbp]
  \centering
  \begin{subfigure}{0.48\textwidth}
    \includegraphics[width=\textwidth]{shape_dark_pt_lambda.pdf}
    \caption{}
  \end{subfigure}
  \hfill
  \begin{subfigure}{0.48\textwidth}
    \includegraphics[width=\textwidth]{shape_dark_decay_len_ctau.pdf}
    \caption{}
  \end{subfigure}
  \caption{The two directions spanned by the training samples of the
  surrogate, shown through generator-level quantities of the truth
  large-radius jets. (a) Dark-hadron transverse momentum across the
  $\mpid/\Lambda_{\mathrm{D}}$ scan at $\ctau = 0.1$~mm. (b)
  Dark-hadron decay length across the lifetime grid of the nominal
  model. Distributions are normalized per logarithmic bin.}
  \label{fig:samplespan}
\end{figure}

\subsection{Detector simulation}
\label{sec:detector}

Events are passed through \textsc{Delphes}~3.5.1
\cite{deFavereau:2013fsa} with the ATLAS detector card. The stock card
propagates charged particles in the solenoidal field and applies
$\pT$- and $\eta$-dependent tracking efficiencies and momentum
resolutions, but leaves the track impact parameters at their generated
values. Since the impact parameters are the observables that carry the
lifetime information, a track-smearing module is inserted after track
reconstruction, carrying the ATLAS transverse and longitudinal
impact-parameter resolutions as a function of track
$\pT$~\cite{ATLAS:IDperf}. The resolution ranges from $80~\mu$m for
tracks below $1$~GeV to $11.5~\mu$m above $30$~GeV in $\dnought$, and
from $140~\mu$m to $47~\mu$m in $\znought$. Every track therefore
carries a smeared $\dnought$ and $\znought$ together with their
uncertainties.

Two simplifications of the tracking model deserve emphasis. The
tracking efficiency in \textsc{Delphes} does not depend on the
production radius of the particle, and no fake tracks are simulated.
Both are appropriate for the lifetime range studied here, in which the
dark-pion decay products originate well inside the innermost pixel
layer, and both would need to be revisited for lifetimes beyond a few
millimeters. Jets are reconstructed from particle-flow objects with
the anti-$k_t$ algorithm~\cite{Cacciari:2008gp,Cacciari:2011ma} with
$R = 0.4$, with the jet energy scale correction of the ATLAS card.

\subsection{Jet reconstruction and object association}
\label{sec:jets}

The same jet recipe is applied on the detector-level and on the
truth-level side, so that the surrogate sees the same jet the tagger
sees. Small-radius anti-$k_t$ $R = 0.4$ jets with $\pT > 20$~GeV and
$|\eta| < 2.5$ are reclustered with the anti-$k_t$ algorithm with
$R = 1.0$ into large-radius jets, following the reclustered-jet
approach used in the ATLAS emerging-jets search~\cite{ATLAS:EJ}. At
detector level the small-radius inputs are the \textsc{Delphes}
particle-flow jets. At truth level they are built from all stable
visible generator particles, excluding neutrinos and muons in the
ATLAS truth-jet convention. Large-radius jets are required to have
$\pT > 200$~GeV and $|\eta| < 2.0$. Truth large-radius jets are kept
down to $150$~GeV so that the surrogate can model the turn-on of the
detector-level threshold, and detector-level and truth jets are matched
one to one within $\dR < 0.5$.

Objects are associated to a large-radius jet through its constituent
small-radius jets: a track or generator particle belongs to a
large-radius jet if it lies within $\dR < 0.4$ of the axis of one of
the small-radius jets that were reclustered into it. This mirrors
ghost association in the experiments and is identical on both sides of
the chain.

\obshead{Detector-level inputs}

The tagger operates on the reconstructed tracks associated to a
large-radius jet, with $\pT > 0.5$~GeV and $|\eta| < 2.5$, ordered by
$\pT$ and truncated at 100 tracks per jet. Each track is described by
its $\pT$, its momentum fraction $\pT/\pT^{\mathrm{jet}}$, its position
relative to the jet axis ($\Delta\eta$, $\Delta\phi$, $\dR$), its
transverse impact parameter signed with respect to the jet axis, its
longitudinal impact parameter, the resolutions $\sigma(\dnought)$ and
$\sigma(\znought)$, its charge, and a categorical flag distinguishing
hadrons, electrons and muons. The jet-signed impact parameter is
positive when the point of closest approach lies on the same side of
the primary vertex as the jet, so that the decay products of a
particle traveling along the jet direction populate the positive tail,
while resolution effects are symmetric.

\obshead{Truth-level inputs}

The surrogate operates on the generator particles associated to a
truth large-radius jet: every stable particle and every decayed
physical particle, that is, Standard Model hadrons and taus as well as
dark hadrons, with $\pT > 0.5$~GeV, ordered by $\pT$ and truncated at
150 particles per jet. The neutrinos are excluded. Each particle is
described by its kinematics relative to the jet, its charge and
status, its production vertex ($L_{xy}$ and $z$), and, for particles
that decay, the decay vertex, the flight length and the direction of
the displacement relative to the momentum, the number of decay
products, and the kinematics of the first two decay products. The PDG
identifiers of the particle and of its first two decay products are
passed through learned embeddings over the vocabulary of identifiers
seen in training, with unseen identifiers mapped to a shared token.
Nothing in this representation refers to a detector. In particular, the
dark pions themselves enter as particles with a displaced decay vertex
and two quark daughters, which is exactly the information a generator
provides for any long-lived state.

Displacements are passed through a symmetric logarithm,
$\mathrm{symlog}_{s}(x) = \mathrm{sign}(x)\ln(1 + |x|/s)$ with a
scale $s$ of $10~\mu\mathrm{m}$ for transverse and $1$~mm for
longitudinal quantities, so that the micrometer-scale impact parameters
of prompt tracks and the millimeter-scale decay lengths of dark pions
occupy comparable dynamic range. Energies and momenta enter
logarithmically, and all continuous inputs are standardized with the
training-set mean and standard deviation. The two input
representations are summarized in Table~\ref{tab:features}.

\begin{table}[htbp]
  \centering
  \small
  \setlength{\tabcolsep}{5pt}
  \begin{tabular}{llc}
    \toprule
    input & description & transform \\
    \midrule
    \multicolumn{3}{l}{\emph{tagger, up to 100 tracks per jet ($\pT > 0.5$~GeV, $|\eta| < 2.5$)}} \\
    $\pT$, $\pT/\pT^{\mathrm{jet}}$ & transverse momentum and jet momentum fraction & $\ln$ \\
    $\Delta\eta$, $\Delta\phi$, $\dR$ & position relative to the jet axis & -- \\
    $\dnought$ & transverse impact parameter, jet-signed & $\mathrm{symlog}_{10\,\mu\mathrm{m}}$ \\
    $\znought$ & longitudinal impact parameter & $\mathrm{symlog}_{1\,\mathrm{mm}}$ \\
    $\sigma(\dnought)$, $\sigma(\znought)$ & impact-parameter resolutions & $\ln$ \\
    $q$ & electric charge & -- \\
    type & hadron, electron or muon & embedded \\
    \midrule
    \multicolumn{3}{l}{\emph{surrogate, up to 150 generator particles per jet ($\pT > 0.5$~GeV)}} \\
    $\pT$, $\pT/\pT^{\mathrm{jet}}$, $E$ & kinematics and jet momentum fraction & $\ln$ \\
    $m$ & mass & -- \\
    $\Delta\eta$, $\Delta\phi$, $\dR$ & position relative to the jet axis & -- \\
    $q$, status & electric charge and generator status & -- \\
    $L_{xy}^{\mathrm{prod}}$, $z^{\mathrm{prod}}$ & production vertex & symlog \\
    $L_{xy}^{\mathrm{dec}}$, $z^{\mathrm{dec}}$, $L^{\mathrm{dec}}$ & decay vertex and flight length & symlog \\
    $\Delta\eta^{\mathrm{dec}}$, $\Delta\phi^{\mathrm{dec}}$ & flight direction relative to the momentum & -- \\
    has decay, $N_{\mathrm{ch}}$ & decay flag and number of decay products & -- \\
    $\pT^{c_{1,2}}$, $\Delta\eta^{c_{1,2}}$, $\Delta\phi^{c_{1,2}}$, $m^{c_{1,2}}$ & kinematics of the first two decay products & $\ln$ ($\pT$) \\
    PDG id, PDG id$^{c_{1,2}}$ & identities of the particle and its decay products & embedded \\
    \bottomrule
  \end{tabular}
  \caption{Input features of the detector-level tagger (reconstructed
  tracks) and of the truth-level surrogate (generator particles), both
  relative to the associated large-radius jet. The symmetric logarithm
  $\mathrm{symlog}_{s}(x) = \mathrm{sign}(x)\ln(1 + |x|/s)$ uses
  $s = 10~\mu\mathrm{m}$ for transverse and $s = 1$~mm for
  longitudinal displacements, and every continuous input is
  standardized after its transform. Categorical inputs are learned
  embeddings over the identifier vocabulary of the training set, with
  unseen identifiers mapped to a shared token.}
  \label{tab:features}
\end{table}

% =====================================================================
\section{The detector-level emerging-jet tagger}
\label{sec:tagger}

\subsection{Architecture and training}

The tagger is a set transformer acting on the tracks of a jet. Each
track is embedded by a two-layer perceptron into a 128-dimensional
representation, a learned class token is prepended, and the set passes
through four pre-normalization transformer encoder layers with eight
attention heads, with the padding slots masked from attention. The
class-token output is normalized and passed through a two-layer head to
a single logit. The model has $0.57$ million parameters. It is trained
with the binary cross-entropy loss using AdamW, a linear warm-up
followed by a cosine learning-rate schedule with peak learning rate
$3\times10^{-4}$, batches of 512 jets, and early stopping on the
validation loss. The six signal lifetimes are pooled, so that a
single tagger serves the full lifetime grid, as in the ATLAS search.

The tagger trains on the training subsamples of the nominal signal
model and of the background of Table~\ref{tab:samples}, with
independent subsamples for validation and for testing. The training
set contains $3.2\times10^6$
large-radius jets, of which $2.0\times10^6$ are signal.

\subsection{Working point and signal region}
\label{sec:sr}

The working point is the logit threshold at which the background jet
efficiency on the test sample is $10^{-3}$, that is, a background
rejection of $10^{3}$. The test background sample of $6.0\times10^5$
jets leaves $600$ jets above threshold, so the working point is a
measurement and not a statistical bound. The signal region requires at
least two large-radius jets with $\pT > 200$~GeV above the working
point. Its efficiency $\effSR$ is the quantity the surrogate is asked
to predict.

Training stops after $21$ epochs, with the best validation loss at
epoch $14$ and a validation accuracy of $0.935$. On the test sample the
area under the ROC curve is $0.979$, and the ROC curves per lifetime
are shown in Fig.~\ref{fig:roc}. The threshold at $10^{3}$ rejection
gives the signal jet efficiencies of Table~\ref{tab:wp}.


\begin{figure}[htbp]
  \centering
  \includegraphics[width=0.6\textwidth]{tagger_roc.pdf}
  \caption{Background rejection as a function of signal jet efficiency
  for the track transformer tagger, evaluated on the test sample, for
  each dark-pion lifetime. The dashed line marks the working point at a
  rejection of $10^{3}$.}
  \label{fig:roc}
\end{figure}

\begin{table}[htbp]
  \centering
  \small
  \setlength{\tabcolsep}{4pt}
  \begin{tabular}{lcccccc}
    \toprule
    $\ctau(\pid)$ & $0.01$~mm & $0.05$~mm & $0.1$~mm & $0.5$~mm & $1$~mm & $5$~mm \\
    \midrule
    jet efficiency at $1/10^{3}$ & $0.250$ & $0.708$ & $0.815$ & $0.909$ & $0.928$ & $0.947$ \\
    signal-region efficiency $\effSR$ & $0.048$ & $0.333$ & $0.431$ & $0.525$ & $0.545$ & $0.550$ \\
    \bottomrule
  \end{tabular}
  \caption{Tagger performance at the working point on the nominal
  signal model, test sample.}
  \label{tab:wp}
\end{table}

The shortest lifetime is the difficult one. At $\ctau = 0.01$~mm the
typical dark-pion flight length is comparable to the impact-parameter
resolution, and the tagger relies on the substructure of the dark
shower rather than on displacement. The tagger does not need to be
optimal for the purpose of this study. It needs to be fixed, so that
the surrogate has a well-defined target, and it needs to be
representative of the class of low-level machine-learned selections
whose reinterpretation is the problem.

% =====================================================================
\section{The truth-level surrogate}
\label{sec:surrogate}

\subsection{Target and training}

The surrogate is a second set transformer with the same architecture
as the tagger, acting on the generator particles of a truth
large-radius jet, with the embedded PDG identifiers concatenated to the
continuous inputs. Its target for a truth jet is the binary outcome of
the full detector-level chain: whether a matched reconstructed jet
exists, passes the kinematic selection and lies above the tagger
working point. Truth jets without a reconstructed partner are
negatives. The surrogate is thus trained to predict the analysis
decision, not to imitate the tagger score, and its output is a
per-jet probability
\begin{equation}
  p_j = P(\text{truth jet } j \text{ is tagged at detector level}),
\end{equation}
obtained from the binary cross-entropy loss, which is a proper scoring
rule and therefore yields calibrated probabilities when the model
capacity and statistics permit.

The surrogate is trained on signal only, on its dedicated training
subsamples of Table~\ref{tab:samples}, which share no event with
anything the tagger has seen, so that the labels are the out-of-sample
decisions of the frozen tagger. The training set combines the nominal model with
the two $\mpid/\Lambda_{\mathrm{D}}$ training points at all six
lifetimes, $3.5\times10^{6}$ truth jets in total, so that the
fragmentation of the dark shower and the amount of displacement vary
independently in training. Background jets are not used, since the
surrogate is never applied to background and jets with labels near
zero only distort the calibration on the signal-like phase space. The
alternative-mass models, the remaining
$\mpid/\Lambda_{\mathrm{D}}$ points, the dark-flavour variants and
the test subsample of the nominal model are held out entirely.

\subsection{From jets to events}

The signal-region efficiency of an event is predicted from the per-jet
probabilities of its truth jets. Treating the jets as independent, the
probability that at least two of $n$ truth jets are tagged is the
Poisson-binomial tail
\begin{equation}
  P_{\geq 2} = 1 - \prod_{j=1}^{n}(1 - p_j)
             - \sum_{k=1}^{n} p_k \prod_{j \neq k}(1 - p_j),
  \label{eq:pb}
\end{equation}
and the predicted efficiency is the mean of $P_{\geq 2}$ over events.
This uses the full probabilistic output of the surrogate. The
alternative of thresholding each $p_j$ at $0.5$ and counting tagged
jets discards the information in the intermediate probabilities and
is reported alongside for comparison. The two agree only when the
surrogate is confident on every jet, and their difference is a
diagnostic of where the truth record fails to determine the detector
response.

% =====================================================================
\section{Results}
\label{sec:results}

\subsection{Closure on the training model}

On the held-out test sample of the nominal model the surrogate
reproduces the signal-region efficiency to $5$--$10\%$ at every
lifetime above $0.01$~mm, with the prediction consistently below the
detector level (ratios $0.90$--$0.95$ in Table~\ref{tab:closure} and
Fig.~\ref{fig:sr}). About half of that offset is the truth-level jet
threshold: a Poisson-binomial reference built from the true
detector-level labels of the same truth jets sits $3$--$5\%$ below the
actual efficiency, because a small fraction of passing reconstructed
jets have no truth partner above $150$~GeV. At $\ctau = 0.01$~mm the
efficiency itself is only $0.05$ and the surrogate overshoots by
$11\%$. In that regime the per-jet probabilities are far from $0$ and
$1$, and thresholding them at $0.5$ predicts an efficiency of $0.003$,
wrong by an order of magnitude, while the Poisson-binomial combination
of Eq.~(\ref{eq:pb}) stays within $11\%$: the probabilistic
combination is not an optional refinement but the difference between a
usable and a useless prediction. The per-jet picture behind these
numbers is in Fig.~\ref{fig:jetpt}, and the calibration of the
probabilities themselves in Fig.~\ref{fig:calib}.

\begin{figure}[htbp]
  \centering
  \begin{subfigure}{0.32\textwidth}
    \includegraphics[width=\textwidth]{sr_efficiency_m5.pdf}
    \caption{}
  \end{subfigure}
  \begin{subfigure}{0.32\textwidth}
    \includegraphics[width=\textwidth]{sr_efficiency_m10.pdf}
    \caption{}
  \end{subfigure}
  \begin{subfigure}{0.32\textwidth}
    \includegraphics[width=\textwidth]{sr_efficiency_m2.pdf}
    \caption{}
  \end{subfigure}
  \caption{Signal-region efficiency as a function of the dark-pion lifetime for (a) the nominal model on the held-out test sample and (b, c) the two dark-pion masses not used in training, $\mpid = 10$ and $2$~GeV. Filled circles are the detector-level efficiency from the reconstructed tracks and the tagger, open squares the surrogate prediction from the truth record alone, Eq.~(\ref{eq:pb}). The lower pads show the ratio of predicted to detector-level efficiency.}
  \label{fig:sr}
\end{figure}

\begin{table}[htbp]
  \centering
  \begin{tabular}{llcccc}
    \toprule
    $\mpid$ & $\ctau(\pid)$ & $\effSR$ actual & $\effSR$ predicted & pred./actual & jet AUC \\
    \midrule
    $5$~GeV & $0.01$~mm & $0.048 \pm 0.001$ & $0.053 \pm 0.001$ & $1.11$ & $0.592$ \\
     & $0.05$~mm & $0.333 \pm 0.002$ & $0.316 \pm 0.001$ & $0.95$ & $0.659$ \\
     & $0.1$~mm & $0.431 \pm 0.002$ & $0.388 \pm 0.001$ & $0.90$ & $0.704$ \\
     & $0.5$~mm & $0.525 \pm 0.002$ & $0.497 \pm 0.002$ & $0.95$ & $0.780$ \\
     & $1$~mm & $0.545 \pm 0.002$ & $0.514 \pm 0.002$ & $0.94$ & $0.805$ \\
     & $5$~mm & $0.550 \pm 0.002$ & $0.524 \pm 0.002$ & $0.95$ & $0.851$ \\
    $10$~GeV & $0.01$~mm & $0.031 \pm 0.001$ & $0.018 \pm 0.001$ & $0.59$ & $0.574$ \\
     & $0.05$~mm & $0.328 \pm 0.002$ & $0.131 \pm 0.001$ & $0.40$ & $0.634$ \\
     & $0.1$~mm & $0.442 \pm 0.002$ & $0.235 \pm 0.001$ & $0.53$ & $0.685$ \\
     & $0.5$~mm & $0.534 \pm 0.002$ & $0.384 \pm 0.002$ & $0.72$ & $0.776$ \\
     & $1$~mm & $0.553 \pm 0.002$ & $0.425 \pm 0.002$ & $0.77$ & $0.806$ \\
     & $5$~mm & $0.570 \pm 0.002$ & $0.492 \pm 0.002$ & $0.86$ & $0.839$ \\
    $2$~GeV & $0.01$~mm & $0.046 \pm 0.001$ & $0.090 \pm 0.001$ & $1.97$ & $0.571$ \\
     & $0.05$~mm & $0.262 \pm 0.002$ & $0.293 \pm 0.001$ & $1.12$ & $0.620$ \\
     & $0.1$~mm & $0.356 \pm 0.002$ & $0.384 \pm 0.002$ & $1.08$ & $0.653$ \\
     & $0.5$~mm & $0.484 \pm 0.002$ & $0.479 \pm 0.002$ & $0.99$ & $0.728$ \\
     & $1$~mm & $0.513 \pm 0.002$ & $0.509 \pm 0.002$ & $0.99$ & $0.759$ \\
     & $5$~mm & $0.483 \pm 0.002$ & $0.486 \pm 0.002$ & $1.01$ & $0.841$ \\
    \bottomrule
  \end{tabular}
  \caption{Signal-region efficiency from the full detector-level chain
  and from the surrogate, with their ratio, and the area under the
  per-jet ROC curve of the surrogate against the detector-level
  decision. The $\mpid = 5$~GeV rows use the held-out test sample of
  the training model. The other masses were never seen in training.}
  \label{tab:closure}
\end{table}

\begin{figure}[htbp]
  \centering
  \begin{subfigure}{0.32\textwidth}
    \includegraphics[width=\textwidth]{jet_eff_vs_pt_m5.pdf}
    \caption{}
  \end{subfigure}
  \begin{subfigure}{0.32\textwidth}
    \includegraphics[width=\textwidth]{jet_eff_vs_pt_m10.pdf}
    \caption{}
  \end{subfigure}
  \begin{subfigure}{0.32\textwidth}
    \includegraphics[width=\textwidth]{jet_eff_vs_pt_m2.pdf}
    \caption{}
  \end{subfigure}
  \caption{Per-jet tagging probability as a function of the truth large-radius jet $\pT$ for (a) $\mpid = 5$~GeV (held-out test sample), (b) $10$~GeV and (c) $2$~GeV, per lifetime. Markers are the detector-level pass fraction, dashed lines the mean surrogate probability in the same bin. The turn-on below $200$~GeV reflects the detector-level jet $\pT$ threshold.}
  \label{fig:jetpt}
\end{figure}

\begin{figure}[htbp]
  \centering
  \begin{subfigure}{0.32\textwidth}
    \includegraphics[width=\textwidth]{calibration_m5.pdf}
    \caption{}
  \end{subfigure}
  \begin{subfigure}{0.32\textwidth}
    \includegraphics[width=\textwidth]{calibration_m10.pdf}
    \caption{}
  \end{subfigure}
  \begin{subfigure}{0.32\textwidth}
    \includegraphics[width=\textwidth]{calibration_m2.pdf}
    \caption{}
  \end{subfigure}
  \caption{Calibration of the surrogate for (a) $\mpid = 5$~GeV (held-out test sample), (b) $10$~GeV and (c) $2$~GeV: the detector-level pass fraction in bins of the predicted probability, with the marker area proportional to the fraction of jets in the bin. A calibrated surrogate lies on the diagonal. No calibration is applied to the surrogate output.}
  \label{fig:calib}
\end{figure}

\subsection{Generalization to unseen dark-pion masses}

The rescaled dark sectors change the multiplicity and the momentum
sharing of the dark hadrons in a jet, and through the dark-pion boost
the distribution of decay lengths at fixed proper lifetime. A surrogate
that had merely memorized the nominal model would fail here in a
characteristic way, predicting the nominal efficiency regardless of
the mass. That is not what happens: the predictions move with the mass,
in the right direction, but not by the right amount. At
$\mpid = 2$~GeV the surrogate closes to better than $12\%$ for every
lifetime above $0.01$~mm. At $\mpid = 10$~GeV it undershoots, by
$14\%$ at $5$~mm and by up to a factor $2.5$ at $0.05$~mm. The error
is in the per-jet probabilities, not in the event-level combination:
the labels-only Poisson-binomial reference closes at every point, so
Eq.~(\ref{eq:pb}) is exact for practical purposes, and the calibration
curves of Fig.~\ref{fig:calib} show the same off-diagonal bands that
the ratios imply.

The direction of the failure is instructive. In the
$\mpid/\Lambda_{\mathrm{D}}$ scan of the training set, fewer and
harder dark hadrons at fixed mass come with a lower tagger efficiency,
and the surrogate has learned that association. The coherently rescaled
$\mpid = 10$~GeV model also produces fewer and harder dark hadrons,
but its tagger efficiency is high, because the displacement per hadron
grows with the boost at fixed proper lifetime. Fragmentation and mass
scale are entangled in the training set with a single dark-pion mass,
and disentangling them requires a second mass in training, the same
lesson at the next level of the parameter hierarchy that the
fragmentation scan taught at the first. The mass points are therefore
the natural next addition to the training mixture, and nothing in the
method limits how many models enter it.

\subsection{Generalization across the dark-sector parameters}

The fragmentation scan is the direct test of the sample-design argument
of Fig.~\ref{fig:samplespan}. A first surrogate trained on the nominal
model alone, with only the lifetime varying, predicted an efficiency
independent of $\mpid/\Lambda_{\mathrm{D}}$, with errors up to a
factor $1.8$ at the ends of the scan: it had learned displacement and
was blind to the structure of the shower. With the two fragmentation
training points added, the prediction follows the measured dependence
across the full scan (Fig.~\ref{fig:scans}a), closing to $4$--$12\%$
at the trained points and at the outermost held-out point, and
undershooting most, by $36\%$, at $\mpid/\Lambda_{\mathrm{D}} = 0.7$,
halfway between two trained points, where the surrogate interpolates a
steeper response than the detector shows. The dark-flavour variants
barely change the detector-level efficiency and the surrogate agrees
(Fig.~\ref{fig:scans}b, ratios $0.87$--$0.88$).

The $Z'$ mass scan probes the momentum dependence with the dark sector
held fixed, on models never used in training. Raising $m_{Z'}$ from
$1.5$ to $3$~TeV hardens the dark-quark spectrum and cuts the
signal-region efficiency by a factor of about $2.5$, mostly through the
per-jet response of the tagger to the changed kinematics. The surrogate
follows it: for every lifetime above $0.01$~mm the ratio of predicted
to detector-level efficiency stays between $0.88$ and $0.98$ at all
three masses (Fig.~\ref{fig:mzp}), the same level of closure as on the
training model itself. The transverse-momentum dependence of the
detector-level selection is learned, not memorized.

\begin{figure}[htbp]
  \centering
  \begin{subfigure}{0.48\textwidth}
    \includegraphics[width=\textwidth]{sr_efficiency_scan_lambda.pdf}
    \caption{}
  \end{subfigure}
  \hfill
  \begin{subfigure}{0.48\textwidth}
    \includegraphics[width=\textwidth]{sr_efficiency_scan_nflav.pdf}
    \caption{}
  \end{subfigure}
  \caption{Signal-region efficiency at $\mpid = 5$~GeV and
  $\ctau = 0.1$~mm as a function of (a) the fragmentation ratio
  $\mpid/\Lambda_{\mathrm{D}}$ and (b) the number of dark-quark
  flavours, detector level versus surrogate.
  Of the scan in (a), the points at $0.2$, $0.5$ and $1.0$ contribute
  to the training mixture through independent samples, and the points
  at $0.35$, $0.7$ and $1.4$ are entirely unseen.}
  \label{fig:scans}
\end{figure}

\begin{figure}[htbp]
  \centering
  \begin{subfigure}{0.32\textwidth}
    \includegraphics[width=\textwidth]{sr_efficiency_mzp2000.pdf}
    \caption{}
  \end{subfigure}
  \begin{subfigure}{0.32\textwidth}
    \includegraphics[width=\textwidth]{sr_efficiency_mzp2500.pdf}
    \caption{}
  \end{subfigure}
  \begin{subfigure}{0.32\textwidth}
    \includegraphics[width=\textwidth]{sr_efficiency_mzp3000.pdf}
    \caption{}
  \end{subfigure}
  \caption{Signal-region efficiency as a function of the dark-pion
  lifetime for mediator masses never used in training, (a)
  $m_{Z'} = 2$~TeV, (b) $2.5$~TeV and (c) $3$~TeV, detector level
  versus surrogate, with the ratio in the lower pads.}
  \label{fig:mzp}
\end{figure}

\subsection{What the surrogate learns}

The scans also expose what the detector-level tagger itself uses,
which is what the surrogate has to reproduce. Across the
fragmentation scan the distribution of the track impact-parameter
significance barely moves, and yet the tagger efficiency at a fixed
number of displaced tracks in the jet shifts by up to $30$ percentage
points between the ends of the scan (Fig.~\ref{fig:effdisp}). The
tagger is not counting displaced tracks: it has learned a joint
response to displacement and to the multiplicity and hardness of the
jet, and any surrogate trained along a single line through that plane
inherits a blind spot. This is the operational argument for training
surrogates on samples that span the physics parameters the analysis
response actually depends on, rather than on one benchmark with a
lifetime grid.

\begin{figure}[htbp]
  \centering
  \includegraphics[width=0.6\textwidth]{reco_eff_vs_n_displaced_lambda.pdf}
  \caption{Detector-level tagger efficiency as a function of the
  number of displaced tracks in the jet ($|\dnought|/\sigma > 3$),
  across the fragmentation scan at $\ctau = 0.1$~mm. At a fixed
  displaced-track multiplicity the efficiency moves by up to $30$
  percentage points, so the tagger uses the structure of the jet
  beyond its displacement.}
  \label{fig:effdisp}
\end{figure}

\stub{Ablation of the truth-level input groups, and the dependence of
the closure on the number of labeled training jets.}

% =====================================================================
\section{Conclusions}
\label{sec:conclusions}

We have built a fully controlled testbed for surrogate-based
reinterpretation of a machine-learned search, from event generation
through an ATLAS-like fast simulation with realistic impact-parameter
resolution, a transformer emerging-jet tagger on reconstructed tracks,
and a two-jet signal region at a working point of $10^{3}$ background
rejection. A second transformer operating only on the generator
particles inside a truth jet, with no reference to any detector
quantity, is trained to predict the tagger decision, and its per-jet
probabilities are combined into event-level signal-region efficiencies
without further approximation.

\todo{Update if numbers move: closure on the training model, generalization
to $\mpid = 10$ and $2$~GeV, calibration.}

The results demonstrate that truth-level particle information is
sufficient to learn a meaningful approximation of the response of a
low-level detector tagger, including at proper lifetimes for which the
displacement is comparable to the impact-parameter resolution and the
tagger relies on substructure. The surrogate is inexpensive in every
relevant sense. It is a single network of under a million parameters,
it consumes the generator record that any theorist can produce, and
it returns a per-event acceptance in milliseconds, so that the
sensitivity of the search to a new model can be estimated without
detector simulation, reconstruction, or access to the
collaboration-internal analysis software.

The controlled setting also makes the limitations explicit. The
surrogate is only as good as the labeled samples it is trained on, and
its accuracy degrades in the directions of parameter space that the
training model does not populate, so the training sample must span the
model variations the reinterpretation is meant to cover. Designing,
training and validating a bespoke surrogate for every analysis would
impose a significant additional burden on analysis teams. A shared
model pre-trained on a broad corpus of simulated events, learning
generic mappings from truth-particle inputs to detector-level outcomes
once and then specialized to individual analyses through lightweight
task-specific heads, is the natural next step, and the scaling of
surrogate accuracy with the number of labeled jets measured here
provides the baseline against which such an approach should be judged.

Although emerging jets serve as the example here, nothing in the
construction is specific to dark showers. The surrogate learns the
detector-level response of a machine-learned tagger from the generator
record of the objects it acts on, and the same recipe applies to any
analysis whose selection is a custom jet tagging algorithm with
constituent-level inputs, whether the discriminating physics is
displacement, substructure, flavor or something else. The extension
from jets to event-level classifiers is direct: the surrogate then
consumes the full generator record of the event and predicts the
event-level decision, and the Poisson-binomial combination is replaced
by the network output itself. The template established here, a
detector-level tagger, a frozen working point, and a truth-level
surrogate validated on models withheld from training, is the material
an analysis team would need to publish for its search to be
reinterpreted, and we encourage its adoption.

\acknowledgments

\todo{acknowledgements}

% =====================================================================
\begin{thebibliography}{99}

\bibitem{Butterworth:2016sqg}
J.~M.~Butterworth, D.~Grellscheid, M.~Kr\"amer, B.~Sarrazin and D.~Yallup,
\emph{Constraining new physics with collider measurements of Standard
Model signatures},
JHEP \textbf{03} (2017) 078, [\texttt{arXiv:1606.05296}].

\bibitem{Buckley:2010ar}
A.~Buckley et al.,
\emph{Rivet user manual},
Comput.\ Phys.\ Commun.\ \textbf{184} (2013) 2803,
[\texttt{arXiv:1003.0694}].

\bibitem{deFavereau:2013fsa}
DELPHES 3 Collaboration, J.~de~Favereau et al.,
\emph{DELPHES 3, A modular framework for fast simulation of a generic
collider experiment},
JHEP \textbf{02} (2014) 057, [\texttt{arXiv:1307.6346}].

\bibitem{Corpe:2025rtz}
L.~D.~Corpe, A.~Haddad and M.~Goodsell,
\emph{Recasting the ATLAS search for displaced hadronic jets in the
ATLAS calorimeter with additional jets or leptons using surrogate
models},
Eur.\ Phys.\ J.\ C \textbf{85} (2025) 1276, [\texttt{arXiv:2502.10231}].

\bibitem{ATLAS:CalRatio}
ATLAS Collaboration,
\emph{Search for neutral long-lived particles that decay into displaced
jets in the ATLAS calorimeter in association with leptons or jets using
$pp$ collisions at $\sqrt{s} = 13$~TeV},
JHEP \textbf{11} (2024) 036, [\texttt{arXiv:2407.09183}].

\bibitem{Schwaller:2015gea}
P.~Schwaller, D.~Stolarski and A.~Weiler,
\emph{Emerging Jets},
JHEP \textbf{05} (2015) 059, [\texttt{arXiv:1502.05409}].

\bibitem{Strassler:2006im}
M.~J.~Strassler and K.~M.~Zurek,
\emph{Echoes of a hidden valley at hadron colliders},
Phys.\ Lett.\ B \textbf{651} (2007) 374, [\texttt{hep-ph/0604261}].

\bibitem{Han:2007ae}
T.~Han, Z.~Si, K.~M.~Zurek and M.~J.~Strassler,
\emph{Phenomenology of hidden valleys at hadron colliders},
JHEP \textbf{07} (2008) 008, [\texttt{arXiv:0712.2041}].

\bibitem{Bai:2013xga}
Y.~Bai and P.~Schwaller,
\emph{Scale of dark QCD},
Phys.\ Rev.\ D \textbf{89} (2014) 063522, [\texttt{arXiv:1306.4676}].

\bibitem{Albouy:2022cin}
G.~Albouy et al.,
\emph{Theory, phenomenology, and experimental avenues for dark showers:
a Snowmass 2021 report},
Eur.\ Phys.\ J.\ C \textbf{82} (2022) 1132, [\texttt{arXiv:2203.09503}].

\bibitem{ATLAS:EJ}
ATLAS Collaboration,
\emph{Search for emerging jets in $pp$ collisions at $\sqrt{s} = 13.6$~TeV
with the ATLAS experiment},
Rept.\ Prog.\ Phys.\ \textbf{88} (2025) 097801,
[\texttt{arXiv:2505.02429}].

\bibitem{Bierlich:2022pfr}
C.~Bierlich et al.,
\emph{A comprehensive guide to the physics and usage of PYTHIA 8.3},
SciPost Phys.\ Codeb.\ \textbf{2022} (2022) 8,
[\texttt{arXiv:2203.11601}].

\bibitem{Carloni:2010tw}
L.~Carloni, J.~Rathsman and T.~Sj\"ostrand,
\emph{Discerning Secluded Sector gauge structures},
JHEP \textbf{04} (2011) 091, [\texttt{arXiv:1102.3795}].

\bibitem{ATLAS:IDperf}
ATLAS Collaboration,
\emph{Track Reconstruction Performance of the ATLAS Inner Detector at
$\sqrt{s} = 13$~TeV},
ATL-PHYS-PUB-2015-018 (2015).
\todo{cross-check the origin of the Delphes ATLAS track-resolution
tables, cited there as arXiv:1405.6569}

\bibitem{Cacciari:2008gp}
M.~Cacciari, G.~P.~Salam and G.~Soyez,
\emph{The anti-$k_t$ jet clustering algorithm},
JHEP \textbf{04} (2008) 063, [\texttt{arXiv:0802.1189}].

\bibitem{Cacciari:2011ma}
M.~Cacciari, G.~P.~Salam and G.~Soyez,
\emph{FastJet User Manual},
Eur.\ Phys.\ J.\ C \textbf{72} (2012) 1896, [\texttt{arXiv:1111.6097}].

\end{thebibliography}

% =====================================================================
\appendix

\section{Pythia Hidden Valley settings}
\label{app:settings}

The signal configuration used for all results in this paper is listed
below. The dark-pion lifetime (\texttt{4900111:tau0}, in mm) and the
dark-sector mass scale are set per grid point at run time. The background samples use the same beam settings with
\texttt{HardQCD:all} and \texttt{PhaseSpace:pTHatMin = 450}.

{\small
\begin{verbatim}
Beams:idA = 2212
Beams:idB = 2212
Beams:eCM = 13600.

! ---- hard process: s-channel Zv ----
HiddenValley:ffbar2Zv = on
4900023:m0 = 1500.
4900023:mWidth = 15.
4900023:onMode = off
4900023:onIfAny = 4900101

! ---- dark sector: SU(3), one light flavor, running coupling ----
HiddenValley:Ngauge = 3
HiddenValley:alphaOrder = 1
HiddenValley:Lambda = 10.
HiddenValley:FSR = on
HiddenValley:fragment = on
HiddenValley:pTminFSR = 11.
HiddenValley:probVector = 0.75

! ---- dark sector masses ----
4900101:m0 = 10.     ! dark quark qv
4900111:m0 = 5.      ! dark pion pi_d
4900113:m0 = 20.     ! dark rho rho_d (rho -> pi pi open)

! ---- decays back to SM ----
4900113:oneChannel = 1 1.0 0 4900111 4900111
4900111:oneChannel = 1 1.0 91 1 -1
4900111:mayDecay = on
\end{verbatim}
}

\section{Detector-simulation settings}
\label{app:delphes}

The \textsc{Delphes} configuration is the ATLAS card distributed with
\textsc{Delphes}~3.5.1 with the following modifications. A
\texttt{TrackSmearing} module is inserted after the track merger and
before the calorimeters, with the $\dnought$, $\znought$, $\phi$ and
$\cot\theta$ resolution tables distributed with \textsc{Delphes} for the
ATLAS detector (\texttt{trackResolutionATLAS}), so that the calorimeter and tau-tagging modules consume the smeared
tracks. Both the particle-flow and the truth jet finders use
$R = 0.4$. The output retains the complete generator record, the
smeared tracks, the particle-flow jets, the truth jets and the missing
transverse momentum.

\end{document}
